\chapter[Introduction (Summary)]{Introduction (Compressed Summary)}

\emph{Based on Lecture 01 -- Introduction. Orientation only: not exam material in itself.}

\section{Exam Essentials}

\begin{keybox}[Exam rules]
\begin{itemize}[leftmargin=*]
  \item Closed book, \textbf{no formula sheet}, basic scientific calculator only (no storage, no NPV functions). On own PC with Respondus; no Excel, no AI.
  \item \textbf{Multiple choice only}: 1 correct answer out of 4, \textbf{no penalty} for wrong answers.
  \item \textbf{Attending} ($\geq 75\%$ attendance): 15 MCQ in 45 minutes (7 Part One + 7 Part Two + 1 Testimonial), worth 0--15/31, plus group homework 0--16/31. Pass if the total is $\geq 18$.
  \item \textbf{Non-attending}: 31 MCQ in 60 minutes (14 Part One + 14 Part Two + 3 Academic Readings), worth 0--31/31.
  \item 31 $\to$ 30 cum laude. \textbf{Discussion Questions are representative of the exam questions.}
\end{itemize}
\end{keybox}

\section{Course Map}

\textbf{Valuation workflow} (each lecture covers one step): Equity story $\to$ Financial statements $\to$ Forecasts $\to$ Cash flows $\to$ Discount rate $\to$ Terminal value $\to$ Multiples $\to$ Value bridge $\to$ Recommendation.

\begin{itemize}
  \item \textbf{Part One} (S.\ Rossi): valuation methods, earnings \& cash flows, business forecasting, special topic on coherent plans, APV \& FCFF, special topic on payout policy. Running case: \textbf{JJ Jones}.
  \item \textbf{Part Two}: cost of capital, growth \& value (D.\ Buosi); DDM, FCFE, multiples (R.\ Zappella); special topics (C.\ Conti).
  \item References: Brealey et al.\ (general), Holthausen--Zmijewski (valuation), McKinsey (practice), Ghezzi--Zanetti (exercises), Kruschwitz--L\"offler (theory).
\end{itemize}

\section{Key Ideas}

\begin{itemize}
  \item \textbf{Five myths}: a valuation is not objective, not exact, not purely quantitative, not timeless; and a more complex model is not a better one (simpler models can do better on average).
  \item \textbf{Why valuation is hard}: we must estimate cash flows, discount rates and risk. The asset-pricing questions are nested: no-arbitrage (least general) $\subset$ portfolio selection $\subset$ equilibrium (most general). Business valuation is a \textbf{general equilibrium pricing problem}.
  \item \textbf{Option view of the firm} ($X$ = asset value, $K$ = face value of debt): equity $=\max(X-K,0)$ = call on the assets $=$ asset $+$ put $-$ riskless bond; risky debt $=\min(X,K)$ = riskless bond $-$ put.
  \item \textbf{Cases}: Ocado IPO (valuation depends on assumptions and perspective; price is negotiated, value is argued; a valuation is tested by new information and is conditional on today's information); Facebook (valuing scalability under extreme uncertainty); Nvidia (if P/E falls while market cap rises, earnings expectations grew faster than price).
\end{itemize}

\begin{keybox}[Takeaway]
Valuation is not clairvoyance: it is a \textbf{disciplined, transparent argument under uncertainty}.
\end{keybox}
